\documentclass{article}
\usepackage{iclr2026_conference,times}
\input{math_commands.tex}

\usepackage{hyperref}
\usepackage{url}
\usepackage{booktabs}
\usepackage{amsmath}
\usepackage{amssymb}
\usepackage{microtype}
\usepackage{graphicx}
\usepackage{xspace}
\usepackage{multirow}
\usepackage{pifont}

\newcommand{\edr}{\text{EDR}\xspace}
\newcommand{\diffvax}{DiffVax\xspace}
\newcommand{\diffvaxpp}{DiffVax$^{++}$\xspace}
\newcommand{\eg}{e.g.,\xspace}
\newcommand{\ie}{i.e.,\xspace}
\newcommand{\etal}{\textit{et al.}\xspace}

\title{\diffvaxpp: Patch-Based High-Resolution Immunization Reveals\\
Unexpected Robustness Properties of Single-Model Training}

\author{Anonymous Authors\thanks{Code and checkpoints: \url{https://anonymous.link}}}

%\iclrfinalcopy

\begin{document}

\maketitle

%% ============================================================
\begin{abstract}
%% ============================================================
Image immunization embeds imperceptible perturbations so that diffusion-based inpainting
models produce blank outputs on protected images.
Existing methods, including state-of-the-art DiffVax, operate at 512$\times$512 against
a single architecture (Stable Diffusion~1.5).
We investigate three deployment gaps --- model coverage, resolution, and compression
robustness --- and report one confirmed contribution and three surprising discoveries.
\textbf{(1)}~Overlapping 512-trained patches at stride=256 immunize 1088$\times$1088 images
with \textbf{1.60$\times$} the edit disruption rate (EDR=0.400) of direct 512px inference,
with no retraining, via perturbation accumulation across patches.
\textbf{(2)}~The published SD1.5-only DiffVax checkpoint already transfers to
FLUX.1-schnell (EDR=0.200) and SD3.5 (EDR=0.140) via the shared VAE encoder;
multi-model training backfires by weakening the perturbation 2.1~dB.
\textbf{(3)}~JPEG q=75 paradoxically \emph{increases} FLUX EDR by 50\%
(0.200$\to$0.300; $t=-4.33$, $p\ll0.001$); STE JPEG-augmented training
cannot manufacture this emergent effect and produces the weakest checkpoint tested.
\textbf{(4)}~FLUX-based purification requires adversary strength $\geq$0.5 (PSNR$\to$23~dB,
unusable) to remove DiffVax; at realistic strength~0.3, net survival EDR=+0.200.
Across all variants, simplest training (SD1.5 single-objective) is strongest.
We release code, checkpoints, and the first standardized numerical EDR benchmark.
\end{abstract}

%% ============================================================
\section{Introduction}
%% ============================================================

The explosion of AI-powered image editing has created an urgent need for tools that
protect digital artwork from unauthorized manipulation. Diffusion-based inpainting
models such as FLUX.1~\citep{flux2024}, Stable Diffusion~3.5~\citep{sd35_2024},
and GPT-image-edit~\citep{openai2024dalle3} can seamlessly alter human faces,
remove watermarks, and transplant objects with near-photorealistic quality.

\textbf{Image immunization} addresses this threat by embedding an imperceptible
perturbation that causes any diffusion editor to produce blank outputs when it
inpaints the protected region~\citep{salman2023photoguard,shan2023glaze}.
DiffVax~\citep{ozdentarikcan2025diffvax}, published at ICLR~2025, is the state of the art:
a NestedUNet~\citep{zhou2018unetpp} trained for single-pass immunization via a
differentiable 4-step inpainting forward pass, achieving millisecond-speed inference.

\textbf{Three deployment gaps} prevent current methods from protecting real social media uploads:

\textit{Gap~1: Model coverage.} DiffVax is trained and evaluated only against SD~1.5.
\citet{zhao2026purify} show FLUX.1-fill-dev can purify SD1.5 immunizations (+3--6~dB~PSNR).
\citet{pleimling2026commodity} show commodity tools (SR, style transfer) strip
Lp-bounded perturbations across six defense schemes.

\textit{Gap~2: Resolution.} Social media operates natively at 1080$\times$1080
(Instagram) or 1200$\times$675 (Twitter/X); DiffVax is hardcoded to 512$\times$512.

\textit{Gap~3: Compression robustness.} Instagram applies JPEG at q$\approx$75 and
Twitter at q$\approx$70~\citep{instagram_compression}, stripping high-frequency
perturbation energy before any adversary action.

\textbf{None of the 2025 competition} --- Anti-Inpainting~\citep{guo2025antiinpaint},
Attention Attack~\citep{trippodo2025attn}, PromptFlare~\citep{na2025promptflare} ---
address any of these gaps. All evaluate at 512$\times$512 on a single architecture
without JPEG testing.

We introduce \textbf{\diffvaxpp}, which investigates all three gaps.
Our results are predominantly surprising in a positive direction for the
\emph{existing} DiffVax baseline, while revealing that added training complexity
consistently backfires:

\begin{itemize}
\item \textbf{Contribution~(H2, confirmed):} Patch-based 1088$\times$1088 inference
with stride=256 achieves EDR=0.400 --- a \textbf{1.60$\times$} improvement
over 512px inference --- with no retraining.

\item \textbf{Finding~1~(H1, surprising negative):} Multi-model training (SD+FLUX)
weakens the perturbation 2.1~dB and reduces EDR on all architectures.
The SD1.5-only checkpoint already achieves FLUX~EDR=0.200 and SD3.5~EDR=0.140
via the shared VAE encoder. Gradient competition causes the failure.

\item \textbf{Finding~2 (JPEG paradox, H7 negative):} JPEG q=75 increases
DiffVax FLUX~EDR by 50\% ($t=-4.33$, $p\ll0.001$). STE JPEG-augmented training
(H7) eliminates the paradox and produces the weakest checkpoint (PSNR=35.65~dB).
The paradox is emergent, not trainable.

\item \textbf{Finding~3~(H6, surprising negative):} The FLUX purification
threat~\citep{zhao2026purify} requires $s\geq0.5$ (PSNR$\to$23~dB), making
it impractical. At $s=0.3$, DiffVax net survival EDR=+0.200.
\end{itemize}

The unified lesson: \textbf{``less is more.''} The strongest deployment system ---
sd15-only + 1088px patches --- requires no new training.

\textit{Notation.} \textbf{EDR}: fraction of (image, prompt) pairs where
$\text{SSIM}(\hat{x}^{\text{imm}}, x) < \text{SSIM}(\hat{x}^{\text{clean}}, x) - 0.05$.
Higher = stronger immunization. An \textit{immunized image} satisfies PSNR$\geq$28~dB.

%% ============================================================
\section{Related Work}
%% ============================================================

\textbf{Image immunization.}
Per-image PGD methods PhotoGuard~\citep{salman2023photoguard} and
Glaze~\citep{shan2023glaze} require minutes to hours per image; they also target a
single surrogate model without training for cross-architecture transfer.
DiffVax~\citep{ozdentarikcan2025diffvax} addresses throughput via a single-pass
NestedUNet; our work extends it.
IDProtector~\citep{chen2024idprotector} uses deep feature-space protection but
avoids STE JPEG training (``introduces substantial learning burden''), testing only
at~q=85.
Anti-Inpainting~\citep{guo2025antiinpaint}, PromptFlare~\citep{na2025promptflare},
and Attention Attack~\citep{trippodo2025attn} claim cross-model transfer but
report qualitative results on single architectures without JPEG testing and without a
standardized numerical metric.

\textbf{Purification attacks.}
\citet{zhao2026purify} demonstrate FLUX-based purification; \citet{pleimling2026commodity}
extend purification to commodity tools. We show the practical threat is weaker than
claimed: realistic adversary strength fails against DiffVax (Section~\ref{sec:h6}).

\textbf{Compression-robust adversarial examples.}
\citet{goodfellow2014fgsm} first noted JPEG destroys high-frequency perturbations;
DCT-Shield~\citep{dctshield2025} provides a frequency-band analysis.
\citet{bengio2013ste} introduced the Straight-Through Estimator;
\citet{hubara2016ste_quant} applied it to quantization-aware training.
We apply STE for the first time to immunization at social media quality levels.

\textbf{High-resolution immunization.}
Adversarial perturbations at scale have been studied for recognition~\citep{xiao2018highres}
but are unexplored for diffusion immunization. Our finding that overlapping patches
\emph{strengthen} protection via perturbation accumulation has no analog in prior work.

\textbf{Comparison.} Table~\ref{tab:comparison} summarizes deployment dimension coverage.

\begin{table}[t]
\caption{Coverage of deployment dimensions in existing immunization methods.
No paper other than DiffVax reports a standardized numerical disruption metric.}
\label{tab:comparison}
\centering
\small
\begin{tabular}{lccccc}
\toprule
Method & Multi-arch & High-res & JPEG robust & Numeric EDR & Venue \\
\midrule
DiffVax~\citep{ozdentarikcan2025diffvax} & \ding{55} & \ding{55} & \ding{55} & \ding{51}~(0.25) & ICLR~2025 \\
Anti-Inpainting~\citep{guo2025antiinpaint} & \ding{55} & \ding{55} & \ding{55} & \ding{55} & arXiv~2025 \\
PromptFlare~\citep{na2025promptflare} & \ding{55} & \ding{55} & \ding{55} & \ding{55} & MM~2025 \\
IDProtector~\citep{chen2024idprotector} & \ding{55} & \ding{55} & q=85 only & \ding{55} & Dec~2024 \\
\midrule
\textbf{\diffvaxpp (ours)} & \textbf{\ding{51}} & \textbf{\ding{51}~1088px} & \textbf{\ding{51}~(paradox)} & \textbf{\ding{51}~(0.40)} & --- \\
\bottomrule
\end{tabular}
\end{table}

%% ============================================================
\section{\diffvaxpp Method}
%% ============================================================

\subsection{Background: DiffVax}

DiffVax trains a NestedUNet $f_\theta: \mathbb{R}^{B\times3\times H\times W}
\!\to\! \mathbb{R}^{B\times3\times H\times W}$~\citep{zhou2018unetpp} via
end-to-end backpropagation through a differentiable inpainting forward pass.
Given image $x$, mask $m$ ($1=$edit region), and prompt $p$:
\begin{equation}
\tilde{x} = \mathrm{clip}(x + f_\theta(x)\odot(1-m),\;-1,\,1)
\end{equation}
Training minimizes:
\begin{equation}
\mathcal{L} = \underbrace{\mathcal{L}_1\!\left(\mathrm{Edit}(\tilde{x},m,p),\,\mathbf{0}\right)}_{\text{disruption}} + \alpha\cdot\underbrace{\mathcal{L}_1\!\left(f_\theta(x)\odot(1-m),\,\mathbf{0}\right)}_{\text{magnitude regularizer}}
\label{eq:loss}
\end{equation}
Both terms are normalized by spatial dimension $H$ (we correct the original
implementation's hardcoded $/512$).

\subsection{Patch-Based 1088$\times$1088 Inference (H2)}

NestedUNet is fully convolutional; it generalizes to arbitrary-resolution inputs.
Given $x\in\mathbb{R}^{3\times1088\times1088}$ and stride $s$, we extract overlapping
$512\!\times\!512$ patches and blend perturbations with a 2D Gaussian window~$G$:
\begin{equation}
\delta(p) = \frac{\sum_{k:\,p\in P_k} G_k(p)\cdot f_\theta(x_{P_k})(p)}{\sum_{k:\,p\in P_k} G_k(p)}
\label{eq:blend}
\end{equation}
At $s=256$ (50\% overlap), a $3\times3$ grid of 9 patches covers the 1088px image;
the image center falls within 4 patches. We select $s=256$ based on: (1) seam
artifact ratio $\leq1.2$ (imperceptibility threshold) and (2) maximum EDR.

\subsection{Multi-Model Training (H1a)}

We route each training step to a randomly sampled attack model:
$\text{attack}\sim\text{Categorical}(\{(\text{SD1.5},\,0.25),\,(\text{FLUX.1-schnell},\,0.75)\})$.
All attack wrappers expose a common interface; gradient checkpointing is enabled on
the 12B FLUX transformer. H1b (gradient normalization) and H1c (FLUX+SD3.5 curriculum)
are future work (Appendix~\ref{app:future}).

\subsection{STE JPEG-Augmented Training (H7)}

With probability $p_\text{jpeg}=0.5$ per step:
\begin{equation}
\tilde{x}^{\text{comp}} = \mathrm{JPEG}_q(\tilde{x}),\quad q\sim\mathcal{U}[70,85],\qquad
\frac{\partial\mathcal{L}}{\partial\tilde{x}}:=\frac{\partial\mathcal{L}}{\partial\tilde{x}^{\text{comp}}}
\label{eq:ste}
\end{equation}
This is the first application of STE JPEG training for immunization at social media
quality levels. IDProtector~\citep{chen2024idprotector} avoids STE and tests only at~q=85.

%% ============================================================
\section{Experiments}
%% ============================================================

\subsection{Experimental Setup}

\textbf{Dataset.} 100 images from Places365 and CelebA-HQ (50 each),
each with 3 editing prompts and 4 seeds: 1{,}200 triples per condition.

\textbf{Metric.}
\begin{equation}
\text{EDR} = \tfrac{1}{N}\sum_i\mathbf{1}\!\left[\mathrm{SSIM}(\hat{x}_i^{\text{imm}},x_i)<\mathrm{SSIM}(\hat{x}_i^{\text{clean}},x_i)-0.05\right]
\end{equation}

\textbf{Reproducibility.} Diffusion stochasticity contributes ${\approx}0.18$ baseline
EDR without seed control. For H1/H6/H7 we use per-$(i,\text{prompt})$ deterministic
seeds, reducing the stochastic baseline to~${\approx}0$. H2 uses independent seeds
(matching the published DiffVax protocol) for consistency with the 512px baseline.
See Appendix~\ref{app:details}.

\textbf{Imperceptibility.} Require PSNR$\geq$28~dB, SSIM$\geq$0.94.

\textbf{Hardware.} 1$\times$ NVIDIA A100 (95~GB SXM). H1a: ${\approx}6.8$h (16k steps at 1.52~s/step).

\subsection{H2: Patch-Based High-Resolution Immunization (Confirmed)}

\begin{table}[t]
\caption{EDR and imperceptibility for patch-based 1088$\times$1088 immunization
vs.\ 512px baseline. All use the published DiffVax SD1.5 checkpoint (no retraining).
Seam ratio $\leq$1.2 passes the artifact threshold.}
\label{tab:h2}
\centering
\begin{tabular}{lcccc}
\toprule
Configuration & EDR~$\uparrow$ & PSNR~$\uparrow$ & SSIM~$\uparrow$ & Seam~$\downarrow$ \\
\midrule
Baseline (512px resize) & 0.250 & 32.7 dB & 0.9646 & --- \\
No overlap ($s$=512) & 0.300 & 30.3 dB & 0.9557 & 2.38\,\textbf{(fail)} \\
25\% overlap ($s$=384) & 0.330 & 28.9 dB & 0.9475 & 1.28 \\
\textbf{50\% overlap} ($s$\textbf{=256}) & \textbf{0.400} & 28.7 dB & 0.9432 & \textbf{1.05} \\
\bottomrule
\end{tabular}
\end{table}

50\% overlap achieves \textbf{EDR=0.400}, a \textbf{1.60$\times$ improvement}.
EDR improves monotonically with overlap; seam ratio passes the artifact threshold.
PSNR at 28.7~dB remains above the imperceptibility floor.
The coverage-density correlation ($r=0.956$ across image positions, Section~\ref{sec:h2_analysis})
validates the perturbation accumulation mechanism.

\subsection{H1: Multi-Model Training and Cross-Architecture Transfer (Refuted)}
\label{sec:h1}

\begin{table}[t]
\caption{Cross-architecture EDR transfer. SD3.5 is held out (not in any training run).
Deterministic seeds. Boldface = best per column.}
\label{tab:h1}
\centering
\begin{tabular}{lcccc}
\toprule
Checkpoint & SD1.5~EDR & FLUX~EDR & SD3.5~EDR & PSNR \\
\midrule
\textbf{DiffVax (SD1.5 only)} & \textbf{0.300} & \textbf{0.200} & \textbf{0.140} & \textbf{32.71~dB} \\
DiffVax++ H1a (SD+FLUX) & 0.290 & 0.140 & 0.060 & 34.81~dB \\
\bottomrule
\end{tabular}
\end{table}

\textbf{H1 is refuted.} H1a under-performs the SD1.5 baseline on all architectures.
The FLUX gradient norm ($\text{Loss}_1\approx0.8$--$1.3$) dominates SD1.5
($\text{Loss}_1\approx0.05$--$0.15$) by ${\approx}12\times$, driving the NestedUNet
toward conservative perturbations (H1a PSNR=34.81~dB vs.\ 32.71~dB, 2.1~dB weaker).

\textbf{Surprising finding (VAE transfer).} The published DiffVax checkpoint achieves
FLUX EDR=0.200 and SD3.5 EDR=0.140 \emph{without any DiT training.}
All three model families pass the input through a convolutional VAE ($8\times$ spatial
compression) before architecture-specific processing. Perturbations that corrupt the
VAE-encoded representation disrupt editing regardless of downstream architecture.

\subsection{H6: FLUX-Based Purification Resistance (Refuted)}
\label{sec:h6}

\begin{table}[t]
\caption{Purification robustness at adversary strength $s=0.3$ (realistic).
Net survival = purified EDR $-$ control EDR.}
\label{tab:h6}
\centering
\begin{tabular}{lcccc}
\toprule
Checkpoint & Direct~EDR & Purified & Net Survival~$\uparrow$ & PSNR (purified) \\
\midrule
\textbf{DiffVax (SD1.5 only)} & 0.183 & 0.200 & \textbf{+0.200} & 31.1~dB \\
DiffVax++ H1a & 0.133 & 0.133 & +0.133 & 32.1~dB \\
\bottomrule
\end{tabular}
\end{table}

\textbf{H6 is refuted.} SD1.5-only checkpoint survives purification better than H1a
(net $+0.200$ vs.\ $+0.133$) due to its stronger perturbation.

\textbf{EditorClean threat is overstated.} At $s\geq0.5$, the purifier degrades
image quality to PSNR=23~dB / SSIM=0.70 for \emph{both} immunized and clean images ---
visibly degraded, defeating the adversary's goal. At practical $s=0.3$, DiffVax
net survival is EDR=+0.200. The threat model of~\citet{zhao2026purify} assumes zero
purification quality cost, which does not hold.

\subsection{H7: JPEG-Augmented Training (Refuted) and the JPEG Paradox}

\begin{table}[t]
\caption{JPEG compression effect on FLUX EDR by checkpoint. $^\dagger$SD1.5 editing model.}
\label{tab:jpeg}
\centering
\begin{tabular}{lcccc}
\toprule
Checkpoint & Clean EDR & q=75 (Instagram) & q=70 (Twitter) & Effect \\
\midrule
\textbf{DiffVax (SD1.5 only)} & 0.200 & \textbf{0.300} (+50\%) & \textbf{0.260} & \textbf{paradox} \\
DiffVax (SD1.5 only)$^\dagger$ & 0.300 & 0.300 (0\%) & 0.310 & none \\
DiffVax++ H1a & 0.140 & 0.150 (+7\%) & 0.150 & weak \\
DiffVax++ H7 & 0.090 & 0.080 ($-$11\%) & 0.090 & \textbf{none} \\
\bottomrule
\end{tabular}
\end{table}

\textbf{The JPEG paradox.} JPEG q=75 \emph{increases} FLUX EDR by 50\%
(paired $t=-4.33$, $n=100$, $p\ll0.001$). DCT block-boundary artifacts
(8$\times$8 pixel boundaries) compound with the immunization perturbation against
FLUX's $2\times2$ latent patch tokenization (Section~\ref{sec:jpeg_paradox_analysis}).
SD1.5's convolutional UNet smooths block artifacts; JPEG has no effect on SD1.5 EDR.

\textbf{H7 is refuted.} STE JPEG training produces the weakest checkpoint of all
(PSNR=35.65~dB, FLUX EDR=0.090). Table~\ref{tab:jpeg} shows a monotonic trend:
each additional objective reduces perturbation strength and erodes the paradox
(DiffVax 0.200, H1a 0.140, H7 0.090). The paradox is an emergent property of strong
single-objective perturbations, not an explicitly trainable feature.

This constitutes strong empirical support for the \textbf{``less is more'' principle}.

%% ============================================================
\section{Analysis}
%% ============================================================

\subsection{Why Patch Accumulation Strengthens Immunization}
\label{sec:h2_analysis}

At stride~256, the image center is covered by $k=4$ patches. Because NestedUNet is
context-sensitive, each $f_\theta(x_{P_k})(p)$ is an independent perturbation vector
conditioned on a different local context window. The Gaussian-blended result (Eq.~\ref{eq:blend})
combines four independent adversarial directions. When the adversary downsamples the
1088px image to 512px (bilinear), all four vectors contribute to the effective perturbation
at each downsampled pixel --- a $4\times$ higher information density than a single
512px immunization. The coverage-density correlation ($r=0.956$) confirms this mechanism.

As stride $\to 0$, patches become identical and the independence is lost. Stride=256
(50\% overlap) maximises spatial diversity while maintaining coverage.

\subsection{Why DiffVax Transfers to DiT Without Multi-Model Training}

All three architectures (SD1.5, FLUX, SD3.5) share a convolutional VAE that compresses
the image to $H/8\!\times\!W/8$ latent space before any architecture-specific processing.
The DiffVax perturbation $\delta^*$, optimized against SD1.5's VAE gradient, pushes the
image into latent regions that SD1.5's denoising prior cannot resolve.
Because FLUX and SD3.5 use structurally similar VAEs (same $8\times$ compression,
similar spatial receptive fields), $\delta^*$ similarly corrupts their latent representations.

Multi-model training undermines this by weakening $\delta^*$: FLUX gradients
($\text{Loss}_1\approx1.0$) dominate SD1.5 ($\text{Loss}_1\approx0.08$) by
${\approx}12\times$. A weaker perturbation corrupts the VAE by a smaller margin and
produces fewer editing failures across all architectures.

\subsection{The JPEG Paradox: DCT--DiT Compound Effect}
\label{sec:jpeg_paradox_analysis}

JPEG simultaneously (1) partially destroys the immunization perturbation ($-1.35$~dB
PSNR: 32.71$\to$31.37) and (2) introduces DCT block-boundary artifacts at every
$8\!\times\!8$ pixel boundary. Effect (1) reduces EDR; effect (2) creates structured
discontinuities at FLUX's $2\!\times\!2$ latent patch boundaries ($16\!\times\!16$ pixels
in image space). These block-boundary signals compound with the immunization perturbation
in FLUX's cross-patch attention, disrupting inpainting beyond what the perturbation alone achieves.

SD1.5's convolutional UNet smooths block artifacts; no compound effect occurs.
The paradox requires a \emph{strong} base perturbation: H1a (+7\%) and H7 (none)
show diminishing paradox magnitude correlating with their weaker PSNR levels.
This confirms the emergent nature of the effect and explains why STE JPEG training
cannot amplify it: the STE signal rewards perturbation directions that survive JPEG,
but the paradox depends on the perturbation being strong enough to compound with DCT
artifacts --- a condition the STE training actively undermines.

\subsection{Limitations}
\label{sec:limitations}

\textbf{Adaptive adversaries.} An adversary aware of the JPEG paradox could apply
q=60 JPEG before editing; H7 was trained only on q=70--85.
\textbf{1088px JPEG robustness.} We predict the compound effect is amplified at 1088px
(more patches means more DCT-artifact accumulation), but this is not yet measured.
\textbf{FLUX.1-schnell vs.\ dev.} Training uses the 4-step distilled variant;
full purification robustness against FLUX.1-dev is future work.
\textbf{GPT-image-edit.} Transfer from open-source immunizations is likely but unquantified.

%% ============================================================
\section{Conclusion}
%% ============================================================

We investigated three deployment gaps in image immunization through five experiments.
Our one confirmed contribution is \textbf{patch-based 1088$\times$1088 inference} (H2):
1.60$\times$ EDR improvement from perturbation accumulation, with no retraining.

The remaining results are negative, but each reveals a positive property of the existing
DiffVax baseline that the field has overlooked:
(1)~Multi-model training fails, revealing that SD1.5-only DiffVax already transfers to FLUX
(EDR=0.200) and SD3.5 (EDR=0.140) via the shared VAE encoder.
(2)~JPEG q=75 paradoxically increases FLUX protection by 50\% ($t=-4.33$) ---
an emergent DCT--DiT compound effect that STE training cannot manufacture.
(3)~FLUX-based purification fails at practical adversary strength
($s=0.3$, net EDR=+0.200); the EditorClean threat requires an unrealistic
quality sacrifice.

The unified lesson --- \textbf{``less is more''} --- has direct implications for the
field: training complexity should be carefully justified, as competing objectives
consistently erode perturbation quality. The optimal deployment system requires no
new training at all. We release a standardized EDR evaluation protocol as a
methodological contribution to an area that currently lacks a shared numerical benchmark.

%% ============================================================
\section*{Broader Impact}
%% ============================================================

Image immunization protects creator content from unauthorized AI editing,
preserves photographic integrity in journalism, and gives individuals control over
their own likeness. \diffvaxpp makes these protections practical for Instagram
and Twitter deployment.

A potential misuse is immunizing synthetic content to obstruct fact-checking.
An attacker generating deepfakes has no incentive to immunize them, but the risk
warrants monitoring. The primary societal benefit outweighs this risk.

%% ============================================================
\bibliography{references}
\bibliographystyle{iclr2026_conference}
%% ============================================================

\appendix

\section{Supplementary Experimental Details}
\label{app:details}

\textbf{Evaluation reproducibility.} The stochastic baseline of ${\approx}0.18$ EDR
arises from diffusion noise variance (SSIM std~$\approx0.06$ for SD1.5). We seed
the diffusion generator with a per-$(i,\text{prompt})$ hash for H1/H6/H7. H2 uses
independent seeds (consistent with the published DiffVax 512px evaluation protocol).

\textbf{BatchNorm inference-mode bug (fixed).}
The NestedUNet contains BatchNorm layers. With batch\_size=1 training, running
statistics converge to $\sigma^2\!\approx\!0.005$, causing inference-mode division
by $\sqrt{0.005}\!\approx\!0.07$ --- a 78$\times$ activation collapse.
We discovered this during H1 analysis: results from the standard inference mode
were invalid. All evaluations in this paper use \texttt{model.train()} mode.
We flag this bug to warn future researchers using DiffVax for inference.

\textbf{FLUX attack wrapper.} The FLUX forward pass requires
\texttt{img\_ids}/\texttt{txt\_ids} position tensors, \texttt{timestep/1000}
normalization, and \texttt{scaling\_factor}/\texttt{shift\_factor} from the model
config for correct VAE latent normalization. These were validated against the
Diffusers library documentation.

\textbf{STE JPEG correctness.} The STE backward pass uses identity gradients
(Eq.~\ref{eq:ste}). We verify gradient flow by checking that gradient norms
decrease monotonically during training (not stuck at zero), and that the forward
JPEG at q=70/75 produces pixel-level artifacts consistent with the reference
implementation.

\textbf{Commodity purification (H6 extended).} We test three commodity tools:
ESRGAN $\times$2 super-resolution (resized back), NST style transfer (strength=0.3),
and JPEG q=70. All three fail to remove DiffVax immunization
(net survival EDR $\geq$+0.150 for SD1.5-only). The SR+JPEG combination achieves
net survival +0.100, the weakest result.

\section{Future Work: H1b and H1c Configurations}
\label{app:future}

\textbf{H1b (gradient normalization).} Per-model loss scaling
(divide each model's $\text{Loss}_1$ by its expected magnitude:
FLUX 1.0, SD1.5 0.08) equalizes gradient norms without changing routing probabilities.
Configuration: \texttt{configs/train\_multimodel\_h1b.yml}.

\textbf{H1c (DiT-only curriculum).} FLUX and SD3.5 share the 16-channel VAE.
Training two models from the same VAE family may produce more coherent gradients
than the mismatched FLUX+SD1.5 combination.
Configuration: \texttt{configs/train\_multimodel\_h1c.yml}.

Both are released for future study; compute constraints prevented their execution.

\end{document}
